\RequirePackage{iftex}
\ifPDFTeX
\documentclass[a4paper,12pt]{article}
\usepackage[hmargin=16mm,vmargin=24mm]{geometry}
\usepackage[T1]{fontenc}
\usepackage[utf8]{inputenc}
\usepackage{CJKutf8}
\else
\documentclass[a4paper,12pt]{jsarticle}
\usepackage[dvipdfmx,hmargin=16mm,vmargin=24mm]{geometry}
\fi
\usepackage{amsmath,amssymb}
\usepackage{enumitem}

\setlength{\parindent}{0pt}
\setlength{\parskip}{0.35em}
\setlist[enumerate]{leftmargin=2.2em,itemsep=0.35em,topsep=0.35em}

\newcommand{\problem}[2][10]{\par\vspace{0.85em}\noindent{\large 問題 #2}\quad（#1点）\quad\ignorespaces}
\newcommand{\R}{\mathbb{R}}

\begin{document}

\ifPDFTeX
\begin{CJK}{UTF8}{min}
\fi

\begin{center}
{\Large 数学1A レポート}\\[0.3em]
各問題10点，満点100点
\end{center}

計算問題では途中式も書くこと。

\problem{1}
関数
\[
f(x,y)=
\begin{cases}
\dfrac{x^2y}{x^2+y^2} & ((x,y)\ne(0,0)),\\[0.7em]
0 & ((x,y)=(0,0))
\end{cases}
\]
について， $f$ は $(0,0)$ で全微分可能かどうか判定せよ。

\problem{2}
関数
\(
f(x,y)=\sqrt{1+x+2y}
\)
を考える。次が成り立つような$A,B$ を具体的に求めよ: 
\[
f(h,k)=1+Ah+Bk+o\left(\sqrt{h^2+k^2}\right)
\qquad ((h,k)\to(0,0))
\]

\problem{3}
曲面
\(
z=x^2+xy+y^2
\)
について，点 $(1,-1,1)$ における接平面の方程式を求めよ。

\problem{4}
ランダウのオー記号について調べ、１００字程度で作文せよ. (何を使ってもよい)

\problem{5}
関数
\(
f(x,y)=\log(1+x+2y)
\)
を考える。点 $(0,0)$ における $2$ 次のテイラーの定理を用いて，
\[
f(h,k)=Ah+Bk+Ch^2+Dhk+Ek^2+o(h^2+k^2)
\qquad ((h,k)\to(0,0))
\]
が成り立つ，$A,B,C,D,E$ を求めよ。

\problem{6}
関数
\(
f(x,y)=x^3+y^3-3xy
\)
の停留点をすべて求め，それぞれが極小点，極大点，鞍点のどれであるか判定せよ。

\problem[20]{7}
平面上の温度を
\(
T(x,y)=100-x^2-2y^2
\)
で表す。点 $(1,2)$ において，次の問いに答えよ。
\begin{enumerate}[label=(\alph*)]
\item $\nabla T(1,2)$ を求めよ。
\item 点 $(1,2)$ を通る等温線 $T(x,y)=T(1,2)$ の接線の方程式を求めよ。
\end{enumerate}

\problem[20]{8}
$a\in\R$ を定数とし，関数
\(
g_a(x,y)=x^4+y^4-4axy
\)
を考える。$a$ の値に応じて，次の問いに答えよ。
\begin{enumerate}[label=(\alph*)]
\item $g_a$ の停留点をすべて求めよ。
\item 各停留点が極小点，極大点，鞍点のどれであるか判定せよ。
\end{enumerate}

\ifPDFTeX
\end{CJK}
\fi

\end{document}
